\section{Probes and charts: the Analysis tab}
\label{sec:analysis}

The strain window has two tabs. \ui{Strain Field} draws one field over the whole region of
interest (Section~\ref{sec:strain}). \ui{Analysis} shows what happens at one place over the
sequence: the displacement of a point, the mean strain of an area, the reading of a virtual
extensometer or a crack-opening gauge, the field along a line, and, with the record of a testing
machine, a stress--strain curve.

On the left of the tab are the placement tools, the reference image with the probes, and a frame
navigator; on the right, the view tabs, the plot controls, the chart and, below it, the probe
table (Figure~\ref{fig:analysis}). A status line at the bottom repeats the last message, which also
goes to the \ui{LOG}. Probes need a finished run, and strain fields need \ui{Compute Strain} on the
\ui{Strain Field} tab; until then the chart says what is missing.

\screenshot{fig-analysis.png}{The \ui{Analysis} tab with three probes on the reference image: a
virtual extensometer (P1), a point between the holes (P2) and a rectangle (P3). The canvas shows
the plotted field at the current frame. The chart shows each probe's reading in the
\ui{Over time} view, the dashed line marking the current frame; the probe table is below
it.}{fig:analysis}

\subsection{Placing a probe}

The toolbar holds seven placement tools as icons; hover over one to see its name and what it
expects. Click a tool, then click on the image.

\begin{center}
\small
\begin{tabularx}{\linewidth}{@{}l l X@{}}
  \toprule
  Tool & Clicks & Result \\
  \midrule
  \ui{Point}     & one & a point (type \ui{Point} in the table) \\
  \ui{Line}      & two: start and end & a line (type \ui{Line}) \\
  \ui{Rectangle} & two: opposite corners & a region (type \ui{Region}) \\
  \ui{Circle}    & two: centre, then the edge & a region \\
  \ui{Polygon}   & one per vertex, double-click to close & a region of three vertices or more \\
  \ui{Virtual extensometer} & two: the gauge points & a line; the chart plots
    \ui{Extensometer strain} \\
  \ui{Crack gauge} & two: one on each side of the crack & a line; the chart plots
    \ui{Crack opening} \\
  \bottomrule
\end{tabularx}
\end{center}

While a tool is armed, a note on the canvas names it and says what it waits for; \kbd{Esc} cancels
the placement. While you draw a line, its length follows the cursor, in pixels or in the length
unit of the physical units, so you can place an extensometer to a prescribed gauge length $L_0$. A
zero-length line, or a shape without area, is ignored.

The extensometer and the crack gauge are ordinary lines. The two tools only choose what the chart
plots once the line is placed, and only if the chart does not already show a reading of the same
kind. Any line can show any gauge reading.

\begin{notebox}
The canvas always shows the reference image (frame~1), whatever frame the chart is on. Probe
coordinates are reference-image pixels and the fields are stored in reference coordinates
(Section~\ref{sec:conventions}), so a probe follows the material it was placed on through the whole
sequence.
\end{notebox}

\ui{Show field} (on by default) colours the reference image with the plotted field at the current
frame, with the field's own colormap and colour range; for a gauge reading, which has no field, it
shows the field of the \ui{Strain Field} tab. The colorbar uses the unit chosen in
\ui{Strain as}. The buttons after the tools fit the image to the view, show it at \ui{100\%} and
zoom in and out; the mouse wheel zooms as well.

\subsection{Selecting, editing and deleting probes}

With no tool armed:

\begin{itemize}
  \item \textbf{Click} a probe to select it: its row is selected and its curve drawn thicker.
    Click empty space, or press \kbd{Esc}, to clear the selection.
  \item \textbf{Drag} a probe to move it, or drag a square \textbf{handle} of the selected probe to
    reshape it: the ends of a line, the corners of a rectangle, the vertices of a polygon, the
    centre (moves) and the right-hand handle (radius) of a circle. A movement of a few screen
    pixels counts as a click. The chart is recomputed when you release the button.
  \item \textbf{Double-click} a probe, or press \kbd{F2}, to rename it. \kbd{Delete} removes the
    selected probe. There is no undo.
  \item \textbf{Drag empty space}, or drag with the middle mouse button, to pan.
\end{itemize}

Markers and labels keep their size on screen at any zoom; the label of a line also gives its
length. New probes are named \file{P1}, \file{P2}, \ldots\ and take the first free colour of an
eight-colour palette. A deleted probe's number is not reused while other probes remain; once every
probe is deleted, the names start again at \file{P1}.

The probe table has the columns \ui{Show} (untick to hide a probe on the canvas and in the chart
without deleting it), \ui{Name} (double-click to edit; it labels the curve and the CSV columns),
\ui{Type} (\ui{Point}, \ui{Line} or \ui{Region}), \ui{Color} and \ui{Note}, which says why a
probe has gaps or is not plotted. Below it, \ui{Color…} recolours the selected probe, \ui{Delete}
removes it, and \ui{Clear All} removes every probe after a confirmation.

\subsection{Choosing what to plot}

A chart plots one \emph{quantity}, and every visible probe that can produce it is drawn: a point, a
line and a region can share a chart of $\varepsilon_{yy}$, but a strain and a displacement never
share an axis.

\begin{description}[leftmargin=0.6cm,style=nextline]
  \item[Plot]
    A field (\ui{Displacement U}, \ui{Displacement V}, \ui{Displacement magnitude} and the ten
    strain components of Section~\ref{sec:strain}) or, below the separator, a gauge reading of a
    line (Table~\ref{tab:gauge-readings}). Until you choose, the chart follows the field of the
    \ui{Strain Field} tab.
  \item[Statistic]
    How the samples of a line or region become one number per frame: \ui{Mean}, \ui{Median},
    \ui{Maximum}, \ui{Minimum}, \ui{Standard deviation}, or \ui{Valid fraction}, the share of the
    probe's material with a valid value on that frame. A point plots its own value for the first
    four and is left out of the last two. Shown for fields, in the \ui{Over time} and
    \ui{Stress–strain} views.
  \item[Strain as] \ui{ratio}, \ui{\%} or \ui{µε}, for strains (the two angles stay in degrees).
  \item[X axis]
    \ui{Frame} (from 1, as in the application); \ui{Time (s)} when physical units with a frame rate
    are on, frame~1 at 0\,s; \ui{Load (N)} and \ui{Stress (MPa)} with load data (see below).
  \item[Min. valid fraction]
    Default 0.50 (0 to 1). A line or region frame is left blank when fewer than this fraction of its
    points are valid, so that a curve does not stay smooth while its sample shrinks to a few points
    next to a crack. Points are not affected.
\end{description}

A point reads the field by bilinear interpolation. A line is sampled once per grid step along its
length. A region reads every grid point inside it, so its statistics are area-weighted; a region
smaller than one grid cell is read at its centre. Displacements and lengths are in pixels, or in
the length unit when physical units are on; $v$ is positive down, as everywhere in RAFTcorr.

\begin{table}[htbp]
\centering
\small
\caption{Gauge readings of a line. $L_0$ is the drawn length, $L$ the distance between the two
ends on the current frame, $\Delta\mathbf{d}$ the displacement of the second end relative to the
first, $\mathbf{e}$ the unit vector from the first end to the second, and $\mathbf{n}$ that vector
turned to the right on screen.}
\label{tab:gauge-readings}
\begin{tabular}{@{}lll@{}}
  \toprule
  Reading & Value & Meaning \\
  \midrule
  \ui{Extensometer strain}      & $(L - L_0)/L_0$ & engineering strain of the gauge \\
  \ui{Extensometer true strain} & $\ln(L/L_0)$ & logarithmic strain \\
  \ui{Elongation ΔL}            & $L - L_0$ & change of the gauge length \\
  \ui{Crack opening}            & $\Delta\mathbf{d}\cdot\mathbf{e}$ & what a clip gauge reads \\
  \ui{Crack sliding}            & $\Delta\mathbf{d}\cdot\mathbf{n}$ & perpendicular to the gauge,
    along the crack \\
  \ui{Crack opening magnitude}  & $|\Delta\mathbf{d}|$ & the whole separation \\
  \bottomrule
\end{tabular}
\end{table}

A gauge reads the displacement at its two ends only, and reads across a crack by design: an
extensometer that spans a crack includes the opening in its elongation. Place a crack gauge with
its ends on the material on each side of the crack, where they stay tracked in every frame.

\subsection{The four views}

\begin{description}[leftmargin=0.6cm,style=nextline]
  \item[Over time]
    Each probe's reading at every frame; a dashed line marks the current frame. Click the chart to
    go to a frame (with a load or stress axis, the nearest one). The frame is shared with the
    \ui{Strain Field} tab and the navigator under the canvas.
  \item[Along the line]
    The field along the selected line, or the newest visible one, on the current frame, against the
    distance from its first end. With \ui{Other frames} (on by default), up to twelve other frames,
    evenly spaced, are drawn in grey behind it. A gauge reading has no profile: choose a field.
  \item[Kymograph]
    The same line on every frame at once: distance against frame (or time), the value as colour in
    the field's colormap and colour range. Load and stress are not offered as the $x$ axis here,
    because they are not evenly spaced.
  \item[Stress–strain]
    Stress, or load without a cross-section, against the plotted quantity, one curve per probe, the
    current frame ringed. With \ui{Extensometer strain} it is a stress--strain curve; with
    \ui{Crack opening}, a load--opening curve. It needs load data.
\end{description}

The small toolbar above the plot pans, zooms and saves the screen picture; while one of its tools
is active, a click does not change the frame.

\subsection{Gaps and marks}

Every frame of every probe is judged on its own, and a gap in a curve is a gap in the data; it is
never bridged. A frame that has a value but lost part of the probe's material (cut out of that
frame's region of interest, or gone out of the image) gets a hollow marker. In the line views,
lost material has its own grey, because it is neither a zero nor a missing measurement. The
legend and the \ui{Note} column say why a probe has gaps:

\begin{center}
\small
\begin{tabularx}{\linewidth}{@{}l X@{}}
  \toprule
  Note & Meaning \\
  \midrule
  \ui{partly untracked from frame N} & part of the material is lost from frame N on; the value
    comes from what remains \\
  \ui{endpoint lost from frame N} & a gauge end left the tracked material \\
  \ui{gaps: too few valid points} & frames below \ui{Min. valid fraction} \\
  \ui{gaps: no strain value} & the virtual strain gauge found too few points (edge, hole,
    trimmed band): use a smaller gauge or a displacement \\
  \ui{not plotted: off the measured area} & the probe lies outside the region of interest \\
  \ui{not plotted: gauges need a line} & a gauge reading for a point or region \\
  \bottomrule
\end{tabularx}
\end{center}

\subsection{Load data from a testing machine}

\ui{Load data…} imports the machine's record as a CSV file (\file{.csv}, \file{.txt},
\file{.dat}). In the \ui{Load Data} dialog, \ui{Choose file…} reads it and shows its first eight
rows. Comma, semicolon and tab separators, decimal commas, a title line and a units row below the
header (whose units join the column names) are all understood. RAFTcorr guesses the columns from
their names; check them:

\begin{description}[leftmargin=0.6cm,style=nextline]
  \item[Load column] The force, and its unit, \ui{N} or \ui{kN}.
  \item[Match rows to frames: By time]
    A \ui{Time column} and an \ui{Offset}, the machine's time at the reference image (if the camera
    started 2\,s after the machine, enter 2\,s). Frame $n$ is at machine time
    $(n-1)/f + \text{offset}$, with $f$ the frame rate of the physical units; the load is
    interpolated there, and frames outside the record get none. Available only with a frame rate.
  \item[By frame number]
    A \ui{Frame column}, and whether \ui{The first image is numbered} 1 or 0. Rows of the same frame
    are averaged.
  \item[Cross-section A0]
    The initial cross-section in mm$^2$, for the engineering stress $F/A_0$ in MPa; 0 plots the load
    only.
\end{description}

A line at the bottom, \ui{Frames with a load: n of N.}, follows every change, so a wrong offset
shows at once. The record is saved with the project. Loading a new image folder clears it; removing
frames keeps it and warns in the \ui{LOG}. \ui{Remove Load Data}, in the same dialog, clears it by
hand.

\subsection{Exporting probes and charts}

The \ui{Export} menu above the chart is available when the chart shows data:

\begin{description}[leftmargin=0.6cm,style=nextline]
  \item[Probe data (CSV)…]
    In \ui{Over time} and \ui{Stress–strain}. One row per frame: \file{frame} (from 1; frame~1 is the reference
    image), \file{image}, \file{time\_s} with a frame rate, \file{load\_N} and \file{stress\_MPa}
    with load data, then three columns per probe: the value (named after probe, field and
    statistic, as \file{P3\_strain\_eyy\_mean}, or probe and reading, as \file{P1\_strain}), its
    \file{\_valid\_fraction} and a \file{\_flag} (\file{ok}, \file{untracked},
    \file{below\_threshold}, \file{unreliable}, \file{endpoint\_lost}, \file{not\_computed},
    \file{no\_data}). An empty cell means no valid measurement, never 0. Strains are written as
    ratios. Header lines starting with \file{\#} give each probe's geometry in reference pixels, the
    units, the frame numbering, the axis convention, the run and strain settings and the source of
    the load. The file is UTF-8 with a byte-order mark, so spreadsheets read non-Latin names.
  \item[Line data (CSV)…]
    In the line views. One row per sample along the line (distance, \file{x\_px}, \file{y\_px}) and
    one column per frame (\file{frame\_1}, \file{frame\_2}, \ldots).
  \item[Chart image…]
    The chart redrawn for a page: white, 16\,$\times$\,10\,cm, 300\,dpi, PNG, SVG or PDF, text
    editable in SVG and PDF, without the frame cursor and the selection.
  \item[Copy chart, Copy plotted data]
    The same figure, or the numbers as drawn and in the units shown as tab-separated text, to the
    clipboard.
\end{description}

Probes with their names, colours and visibility, and the load record, are saved with the project
(Section~\ref{sec:projects}); the chart settings are not.
